\clearpage
\begin{examplebox}
\textbf{例 1：先检查三角形，再按比例分配。}
在 $\triangle ABC$ 中，$AB=9,AC=6,BC=10$，$AD$ 是内角平分线，求 $BD,DC$。
三边满足 $|9-6|<10<9+6$，三角形存在。
由 $BD:DC=9:6=3:2$，设 $BD=3k,DC=2k$，则
\[
 3k+2k=10\quad\Longrightarrow\quad k=2.
\]
所以 $\boxed{BD=6,\ DC=4}$。
\textbf{回验：} $6+4=10$，$6/4=9/6$。
\begin{center}
\begin{tikzpicture}[font=\small]
\draw[black!30,rounded corners=2pt] (0,0) rectangle (14.4,8.2);
\node[anchor=west,font=\small\bfseries] at (0.5,7.72) {例 1：同一长度单位下，正常三角形与退化情形对照};
% MODEL 9 6 10
\begin{scope}[shift={(4.2,4.4)},scale=0.4]
\coordinate (A) at (7.25,5.3326822519);
\coordinate (B) at (0,0);
\coordinate (C) at (10,0);
\coordinate (D) at (6,0);
\coordinate (A0) at (9,-6.75);
\coordinate (B0) at (0,-6.75);
\coordinate (C0) at (15,-6.75);
\draw[blue,thick] (A)--(B);\draw[teal!80!black,thick] (A)--(C);\draw[thick] (B)--(C);\draw[orange!85!black,very thick] (A)--(D);
\fill (A) circle (1.3pt) node[above=5pt] {$A$};
\fill (B) circle (1.3pt) node[below left=5pt] {$B$};
\fill (C) circle (1.3pt) node[below right=5pt] {$C$};
\fill (D) circle (1.3pt) node[below=5pt] {$D$};
\path (A)--(B) node[midway,above left=4pt] {$AB=9$};
\path (A)--(C) node[midway,above right=4pt] {$AC=6$};
\draw[red!75!black,thick] (B0)--(C0);
\fill[red!75!black] (A0) circle (2pt) node[below=6pt] {$A$};
\fill[red!75!black] (B0) circle (2pt) node[below=6pt] {$B$};
\fill[red!75!black] (C0) circle (2pt) node[below=6pt] {$C$};
\end{scope}
\draw[blue,|-|] ($(B)+(0,-0.85)$)--($(D)+(0,-0.85)$) node[midway,below=3pt] {$BD=6$};
\draw[teal!80!black,|-|] ($(D)+(0,-0.85)$)--($(C)+(0,-0.85)$) node[midway,below=3pt] {$DC=4$};
\node[anchor=west] at (.7,6.3) {修正后 $BC=10$};
\node[anchor=west,text=red!75!black] at (.7,2.2) {原题 $BC=15$};
\node[anchor=west,text=red!75!black] at (.7,.8) {$9+6=15$：三点共线};

\end{tikzpicture}
\end{center}

\textbf{原稿勘误：} 原题把 $BC$ 写成 $15$，但 $9+6=15$，三点只能共线。
形式上列出比例方程不等于几何对象存在；本题已把 $BC$ 修正为 $10$，图中上下两组使用相同长度单位。

\tcbbreak
\textbf{例 2：逆用比例证明角平分线。}
在非退化 $\triangle ABC$ 中，$AB=15,AC=9$，$D$ 在 $BC$ 内部，且 $BD=10,DC=6$。
求 $BC$ 并证明 $AD$ 平分 $\angle BAC$。
\begin{center}
\begin{tikzpicture}[font=\small]
\draw[black!30,rounded corners=2pt] (0,0) rectangle (14.4,8.2);
\node[anchor=west,font=\small\bfseries] at (0.5,7.72) {例 2：只给长度，先验证比例，再证明两个角相等};
% MODEL 15 9 16
\begin{scope}[shift={(2.4,2.2)},scale=0.6]
\coordinate (A) at (12.5,8.2915619759);
\coordinate (B) at (0,0);
\coordinate (C) at (16,0);
\coordinate (D) at (10,0);
\draw[blue,thick] (A)--(B);\draw[teal!80!black,thick] (A)--(C);\draw[thick] (B)--(C);\draw[orange!85!black,very thick] (A)--(D);
\fill (A) circle (1.3pt) node[above=5pt] {$A$};
\fill (B) circle (1.3pt) node[below left=5pt] {$B$};
\fill (C) circle (1.3pt) node[below right=5pt] {$C$};
\fill (D) circle (1.3pt) node[below=5pt] {$D$};
\path (A)--(B) node[midway,above left=4pt] {$AB=15$};
\path (A)--(C) node[midway,above right=4pt] {$AC=9$};
\end{scope}
\draw[blue,|-|] ($(B)+(0,-1)$)--($(D)+(0,-1)$) node[midway,below=3pt] {$BD=10$};
\draw[teal!80!black,|-|] ($(D)+(0,-1)$)--($(C)+(0,-1)$) node[midway,below=3pt] {$DC=6$};

\end{tikzpicture}
\end{center}

先看位置：$D$ 在 $B,C$ 之间，所以 $BC=10+6=\boxed{16}$。
再核比例：
\[
 \frac{BD}{DC}=\frac{10}{6}=\frac53=\frac{15}{9}=\frac{AB}{AC}.
\]
由逆定理，$\boxed{\angle BAD=\angle DAC}$。
图中没有预先标等角，等角是推理所得，不能由画图先假定。
\textbf{回验：} 三边满足 $|15-9|<16<15+9$，分点公式给出 $BD=16\times15/24=10$。

\textbf{同图变式：} 若改给 $AD$ 已是角平分线、$BD-DC=4$，仍有 $BD:DC=5:3$。
设 $BD=5t,DC=3t$，由 $2t=4$ 得 $t=2$，故 $BC=8t=16$。
这里“用差”来自额外已知条件，不表示内部分点满足外分差关系。

\tcbbreak
\textbf{例 3：直角三角形，先找齐两邻边。}
在 $\triangle ABC$ 中，$\angle B=90^\circ$，$AB=5,BC=12$，$AD$ 是内角平分线，求 $BD$。
\begin{center}
\begin{tikzpicture}[font=\small]
\draw[black!30,rounded corners=2pt] (0,0) rectangle (14.4,8.2);
\node[anchor=west,font=\small\bfseries] at (0.5,7.72) {例 3：先用勾股求斜边，再按 $5:13$ 分配 $BC$};
% MODEL 5 13 12
\begin{scope}[shift={(2.4,2.2)},scale=0.8]
\coordinate (A) at (0,5);
\coordinate (B) at (0,0);
\coordinate (C) at (12,0);
\coordinate (D) at (3.3333333333,0);
\draw[blue,thick] (A)--(B);\draw[teal!80!black,thick] (A)--(C);\draw[thick] (B)--(C);\draw[orange!85!black,very thick] (A)--(D);
\fill (A) circle (1.3pt) node[above=5pt] {$A$};
\fill (B) circle (1.3pt) node[below left=5pt] {$B$};
\fill (C) circle (1.3pt) node[below right=5pt] {$C$};
\fill (D) circle (1.3pt) node[below=5pt] {$D$};
\path (A)--(B) node[midway,left=5pt] {$AB=5$};
\path (A)--(C) node[midway,above right=4pt] {$AC=13$};
\pic[draw,angle radius=3mm] {right angle=A--B--C};
\end{scope}
\draw[blue,|-|] ($(B)+(0,-1)$)--($(D)+(0,-1)$) node[midway,below=3pt] {$10/3$};
\draw[teal!80!black,|-|] ($(D)+(0,-1)$)--($(C)+(0,-1)$) node[midway,below=3pt] {$26/3$};

\end{tikzpicture}
\end{center}

勾股定理给出 $AC=\sqrt{5^2+12^2}=13$。
角 $A$ 的邻边是 $AB,AC$，所以分段比为 $5:13$，不是 $5:12$。
\[
 BD=\frac{BC\cdot AB}{AB+AC}
    =\frac{12\times5}{5+13}
    =\boxed{\frac{10}{3}}.
\]
\textbf{回验：} $DC=12-10/3=26/3>0$，且 $(10/3)/(26/3)=5/13$。
\end{examplebox}
